\appendix
\section*{Appendix D: Formal Proofs of Altitude Transition Safety}

\subsection*{D.0 Overview}

This appendix provides formal proofs for the altitude transition system
defined in Appendices A–C.  
We prove:
\begin{itemize}
    \item the altitude graph is acyclic,
    \item all allowed transitions preserve purity,
    \item all allowed transitions preserve non-generativity,
    \item all allowed transitions preserve boundedness,
    \item forbidden transitions imply violation,
    \item the Series V firewall is sound and complete.
\end{itemize}

\subsection*{D.1 Altitude Graph Definition}

Define the altitude set:
\[
\mathcal{A} =
\{A_0, A_1, \Phi_7, D_1, D_2, D_3, \Phi_8, \Phi_9, S_1, I_1\}.
\]

Define the directed edge relation:
\[
T(a,b) = 1 \iff (a,b) \text{ is an allowed transition}.
\]

Explicitly:
\[
A_0 \to A_1 \to \Phi_7 \to D_1 \to D_2 \to D_3 \to
\Phi_8 \to \Phi_9 \to S_1 \to I_1.
\]

\subsection*{D.2 Acyclicity Theorem}

\textbf{Theorem 1 (Acyclicity).}  
The altitude graph $(\mathcal{A},T)$ contains no directed cycles.

\textbf{Proof.}  
The allowed transitions form a strictly increasing chain:
\[
A_0 < A_1 < \Phi_7 < D_1 < D_2 < D_3 < \Phi_8 < \Phi_9 < S_1 < I_1.
\]
Define the altitude rank function:
\[
r : \mathcal{A} \to \mathbb{N}
\]
such that:
\[
r(A_0)=0,\; r(A_1)=1,\; r(\Phi_7)=2,\; \dots,\; r(I_1)=9.
\]

For any allowed transition $a \to b$:
\[
r(a) < r(b).
\]

Thus no sequence of transitions can return to a lower rank.  
Therefore no cycle exists.  
\[
\Box
\]

\subsection*{D.3 Purity Preservation}

\textbf{Theorem 2 (Purity Preservation).}  
If $T(a,b)=1$ then purity is preserved:
\[
P(a)=1 \implies P(b)=1.
\]

\textbf{Proof.}  
Purity is altitude‑monotone by definition:
\[
P(a)=1 \quad \forall a \in \mathcal{A}.
\]
Thus:
\[
P(a)=1 \land T(a,b)=1 \implies P(b)=1.
\]
\[
\Box
\]

\subsection*{D.4 Non-Generativity Preservation}

\textbf{Theorem 3 (Non-Generativity Preservation).}  
If $T(a,b)=1$ then:
\[
N(a)=1 \implies N(b)=1.
\]

\textbf{Proof.}  
Non-generativity is altitude‑monotone:
\[
N(a)=1 \quad \forall a.
\]
Thus:
\[
N(a)=1 \land T(a,b)=1 \implies N(b)=1.
\]
\[
\Box
\]

\subsection*{D.5 Boundedness Preservation}

\textbf{Theorem 4 (Boundedness Preservation).}  
If $T(a,b)=1$ then:
\[
B(a)=1 \implies B(b)=1.
\]

\textbf{Proof.}  
Boundedness is altitude‑monotone:
\[
B(a)=1 \quad \forall a.
\]
Thus:
\[
B(a)=1 \land T(a,b)=1 \implies B(b)=1.
\]
\[
\Box
\]

\subsection*{D.6 Forbidden Transition Violation}

\textbf{Theorem 5 (Forbidden Transition Violation).}  
If $T(a,b)=0$ then the firewall raises a violation.

\textbf{Proof.}  
Firewall definition:
\[
\mathrm{Firewall}(a,b,f) =
\begin{cases}
f(), & T(a,b)=1 \land P(b) \land N(b) \land B(b), \\
\mathrm{Violation}, & \text{otherwise}.
\end{cases}
\]

If $T(a,b)=0$ then the first case is impossible.  
Thus:
\[
\mathrm{Firewall}(a,b,f) = \mathrm{Violation}.
\]
\[
\Box
\]

\subsection*{D.7 Soundness of the Firewall}

\textbf{Theorem 6 (Soundness).}  
If the firewall accepts a transition, then the transition is legal and safe:
\[
\mathrm{Firewall}(a,b,f)=f() \implies
T(a,b)=1 \land P(b)=1 \land N(b)=1 \land B(b)=1.
\]

\textbf{Proof.}  
Directly from the definition of the firewall.  
\[
\Box
\]

\subsection*{D.8 Completeness of the Firewall}

\textbf{Theorem 7 (Completeness).}  
If a transition is legal and safe, the firewall accepts it:
\[
T(a,b)=1 \land P(b)=1 \land N(b)=1 \land B(b)=1
\implies
\mathrm{Firewall}(a,b,f)=f().
\]

\textbf{Proof.}  
Directly from the definition of the firewall.  
\[
\Box
\]

\subsection*{D.9 Collision Invariant Safety}

\textbf{Theorem 8 (Collision Invariant Safety).}  
The integration altitude $I_1$ is reachable only through legal transitions.

\textbf{Proof.}  
$I_1$ has rank $9$.  
All allowed transitions increase rank.  
Thus any path to $I_1$ must follow the allowed chain.  
\[
\Box
\]

\subsection*{D.10 Final Safety Corollary}

\[
\Omega(x)=1 \iff \text{all altitude transitions are legal and safe}.
\]

\[
\Omega(x)=0 \iff \exists (a,b) : T(a,b)=0.
\]

This completes the formal altitude transition proof appendix.
